\section{Introduction}
\label{sec:intro}

An optimizing compiler may replace a program with any program that has the same
observable behavior, commonly referred to as the as-if rule of ISO
C++~\cite{n4950}. That rule lets compilers reorder, inline, vectorize,
and eliminate code, and it makes the source the complete specification. Every
detail the source fixes is one the compiler must preserve, whether the developer
cared about it or not.

Source code mixes two kinds of decisions. Domain invariants are what the program
must actually guarantee, such as packets leaving a queue in the order they
arrived and the queue holding at least a thousand packets before it drops.
Incidental choices are details the code had to fix in order to run, such as a
capacity of exactly 1000 or a remainder instruction that maps a counter onto
that capacity. A packet queue written for bursts of 800 packets with 25\%
headroom needs at least 1000 slots. The developer writes 1000, and each enqueue
computes \texttt{slot = i \% 1000}, which compiles to a chain of multiplications
and shifts. With 1024 slots the same index is a bitwise \texttt{and}. Both sizes
meet the requirement, and the as-if rule allows only the size in the source,
because the two queues disagree once they hold 1000 packets.

Developers split these decisions by hand. Folly's single-producer queue rounds
its backing array up to a power of two and indexes with a mask, while
\texttt{isFull} still respects the capacity the caller asked
for~\cite{folly_pcq}. DPDK's exact-size rings do the same
separation~\cite{dpdk_rte_ring}. ClickHouse sizes hash tables in powers of two
and places keys with a mask~\cite{clickhouse_hashtable}. Those libraries write
the split into the source, and a compiler bound by as-if cannot introduce it from
a baseline that used 1000 and a remainder. This paper asks for that compiler.

Compiling against a specification instead of the source is an old idea. Leroy
defines compilers that preserve a specification rather than exact
behavior~\cite{leroy2009backend}. Refinement calculi derive implementations from
specifications~\cite{back1998refinement,morgan1994programming}. Fiat, Cozy, and
Primrose choose among the implementations a specification admits, Primrose
as annotations on Rust~\cite{delaware2015fiat,loncaric2016cozy,Qin2023Primrose}. STOKE
optimizes floating-point code against a stated error bound rather than
equivalence~\cite{schkufza2014fp}. This paper asks what that idea is, as a
compilation problem for ordinary systems code, and what an automated compiler
can actually do with it.

The answer has three layers.

\paragraph{The Ideal Problem.}
The developer states a functional contract \(\Phi_{\text{func}}\) over
observable behavior, including named implementation parameters such as a
capacity of at least 1000, together with a resource budget and a cost
objective. Programs that satisfy \(\Phi_{\text{func}}\) are legal, and
those that also satisfy the budget are feasible. The ideal compiler
returns a cheapest feasible legal program. That legal set contains every
as-if variant of a correct baseline, and more programs exactly where
\(\Phi_{\text{func}}\) leaves a decision open.

\paragraph{A Restricted Compiler.}
\textsc{IntentCompile} proposes candidates from a generator (rewrites and
parameter holes), admits those a sound verifier proves within a time budget,
and hands the winner to an unmodified production compiler. What it can use
is the verified fragment of the generated set, and it emits \(P_1\) or a
program cheaper on the supplied \(\mathcal{O}\).

\paragraph{An Instance.}
\textsc{IntC} is a C instance of \textsc{IntentCompile}. It takes a C
baseline, a contract written as parameter constraints plus CBMC harnesses,
and an optional benchmark.
It rewrites remainder, division, and multiplication by a parameter. It also
applies representation rewrites: pad-and-mask, which grows a private store
while keeping API capacity, and structure-of-arrays. Arithmetic rewrites drop
a sqrt from a disk test or a double accumulator from an inner product. It fills
parameter holes by search near the baseline, admits candidates a bounded model
checker proves, and selects by measurement. On the FIFO ring, that winner is
the 1024-cell mask, \(2.3\times\) on the named benchmark. What it claims of an admitted
program is \(\models_K\): the harness holds on every execution the bound
covers. It emits reviewable source so a developer can inspect an admitted
change as C, including a change that disagrees with the baseline on an
observable result.

\paragraph{Contributions.}
\begin{itemize}
  \item The compilation problem: contracts with named implementation
  parameters, four kinds of slack, and an ideal objective over the legal set
  (\S\ref{sec:ideal}).
  \item A restricted compiler whose search space is a generator, whose
  admission is a sound verifier, including the bounded judgment
  \(\models_K\), and whose backend is an unmodified production compiler
  (\S\ref{sec:compiler}).
  \item The \textsc{IntentCompile} pipeline and \textsc{IntC}, a C instance
  of it (\S\ref{sec:pipeline}, \S\ref{sec:implementation}).
  \item Evaluation of \textsc{IntC} (\S\ref{sec:evaluation}): whether it
  runs on extracts from real-world projects, whether as-if compilation and
  the internals behave as claimed, and whether ranking under the same
  \texttt{/O2} emits a cheaper program on the named driver or keeps
  \(P_1\). On the ring that driver is \(2.3\times\).
\end{itemize}

Section~\ref{sec:overview} walks through the ring buffer and what existing
tools can do with it. Sections~\ref{sec:ideal}--\ref{sec:compiler} are the two
theory layers. Sections~\ref{sec:pipeline}--\ref{sec:implementation} are the
algorithm and the artifact. Section~\ref{sec:evaluation} answers the three questions of \textsc{IntC}.
Section~\ref{sec:related} places the work. Section~\ref{sec:discussion}
names the risks that remain. Section~\ref{sec:conclusion} closes.
